\section{隐私保护与精准遗忘}\label{sec:forget}

\subsection{敏感信息的自动保护}

记忆模块内置六类敏感信息检测：电子邮箱、手机号、身份证号（含校验码验证）、银行卡号
（含 Luhn 验证）、API 密钥与口令字段。凡是要写入记忆的内容，无论来自对话、工具结果
还是配置，都会先经过检测；命中部分在入库前就被替换为 \code{[REDACTED:类型]} 占位符，
原始值不落盘。这项保护在服务每次启动时还会对已有记忆文件做一次全量复扫，防止历史
数据中残留敏感内容。你可以让助手「扫描一下记忆里的敏感信息」主动查看当前的保护状态。

\subsection{用一句自然语言遗忘}

精准遗忘是你的权利，也是一等公民功能。直接用自然语言向助手下达遗忘指令即可，支持
两类修饰语法让指令足够精确：

\begin{lstlisting}[language={}]
忘掉我的所有住址
忘掉我的所有住址，但保留城市偏好          # 排除语法
只忘掉来源是手动配置的偏好                # 来源限定语法
\end{lstlisting}

遗忘永远不会因为一句话就立即执行。助手首先解析你的意图，生成一份\textbf{删除预案}：
列出将要删除的每条内容（以脱敏预览的形式）、保留的内容，以及预案摘要。预案在
五分钟内有效，你需要明确确认后才会执行。确认时系统校验预案未被篡改，随后物理删除
目标数据及其关联，并\textbf{重新读取磁盘验证已无任何残留}，最后在审计记录中留下这次
操作的事实（记录本身不含被删内容）。涉及删除的确认与工作流模板的激活一样，在权限
层被强制要求人工审批，代理无法替你越过这一步。

如果指令含糊到无法确定删除范围（例如「忘掉一些东西」），助手会拒绝执行并请你说明
具体目标——这是刻意设计的安全行为。目前可被遗忘指令识别的记忆对象覆盖常见类型
（住址、城市、电话、邮箱、身份证、偏好等）\todo{遗忘意图可识别范围随实现扩展后
更新本清单}。

\subsection{数据的所有权与彻底清除}

所有记忆都是你工作区里的普通文件，拥有最终控制权：删除 \code{preferences.json} 即
清空全部偏好，删除 \code{structured_memory.json} 即清空知识库，删除 \code{rag_index.json}
即清空语义索引。应用卸载不会遗留工作区之外的记忆数据。
